\chapter{2014年大纲卷（文）}

\section{单选题}

\begin{problem}
设集合 \(M = \{1, 2, 4, 6, 8\}\)，\(N = \{1, 2, 3, 5, 6, 7\}\)，则 \(M \cap N\) 中元素的个数为（\quad）

\choices
    {\(2\)}
    {\(3\)}
    {\(5\)}
    {\(7\)}
\end{problem}

\begin{answer}
B
\end{answer}

\begin{solution}
由 \(M\) 与 \(N\) 的定义，\(M\cap N=\{1,2,6\}\)，所以其中有 \(3\) 个元素，故选 B.
\end{solution}

\begin{problem}
已知角 \(\alpha\) 的终边经过点 \((-4,3)\)，则 \(\cos \alpha =\)（\quad）

\choices
    {\( \frac{4}{5} \)}
    {\( \frac{3}{5} \)}
    {\(-\frac{3}{5}\)}
    {\(-\frac{4}{5}\)}
\end{problem}

\begin{answer}
D
\end{answer}

\begin{solution}
点 \((-4,3)\) 到原点的距离为 \(5\)，而终边经过第二象限，故
\(\cos\alpha=\frac{-4}{5}\)，故选 D.
\end{solution}

\begin{problem}
不等式组 \(
\begin{cases}
x(x + 2) > 0, \\
|x| < 1
\end{cases}
\) 的解集为（\quad）

\choices
    {\( \{x \mid -2 < x < -1\} \)}
    {\(\{x\mid -1 < x < 0\}\)}
    {\( \{x \mid 0 < x < 1\} \)}
    {\( \{x \mid x > 1\} \)}
\end{problem}

\begin{answer}
C
\end{answer}

\begin{solution}
由 \(x(x+2)>0\) 得 \(x<-2\) 或 \(x>0\)，由 \(|x|<1\) 得 \(-1<x<1\)，两者同时成立时 \(0<x<1\)，故选 C.
\end{solution}

\begin{problem}
已知正四面体 \(ABCD\) 中，\(E\) 是 \(AB\) 的中点，则异面直线 \(CE\) 与 \(BD\) 所成角的余弦值为（\quad）

\choices
    {\(\frac{1}{6}\)}
    {\(\frac{\sqrt{3}}{6}\)}
    {\(\frac{1}{3}\)}
    {\(\frac{\sqrt{3}}{3}\)}
\end{problem}

\begin{answer}
B
\end{answer}

\begin{solution}
设正四面体的棱长为 \(s\)，建立直角坐标系，取
\(A=(0,0,0)\)，\(B=(s,0,0)\)，\(C=(\frac{s}{2},\frac{\sqrt{3}s}{2},0)\)，
\(D=(\frac{s}{2},\frac{\sqrt{3}s}{6},\frac{\sqrt{6}s}{3})\)，则
\(E=(\frac{s}{2},0,0)\). 取
\(\overrightarrow{EC}\) 与 \(\overrightarrow{BD}\) 为两异面直线的方向向量，有
\(\overrightarrow{EC}=(0,\frac{\sqrt{3}s}{2},0)\)，
\(\overrightarrow{BD}=(-\frac{s}{2},\frac{\sqrt{3}s}{6},\frac{\sqrt{6}s}{3})\)，因此
\[
\cos\theta=\frac{|\overrightarrow{EC}\cdot\overrightarrow{BD}|}{|\overrightarrow{EC}|\,|\overrightarrow{BD}|}
=\frac{\frac{s^{2}}{4}}{\frac{\sqrt{3}s}{2}\cdot s}=\frac{\sqrt{3}}{6}
\]
故选 B.
\end{solution}

\begin{problem}
函数 \( y = \ln (\sqrt[3]{x} + 1) \) （\( x > -1 \)） 的反函数是（\quad）

\choices
    {\( y = (1 - \e^{x})^{3}  \)（\(x > -1\)）}
    {\( y = (\e^{x} - 1)^{3}  \)（\(x > -1\)）}
    {\( y = (1 - \e^{x})^{3}  \)（\(x \in \R\)）}
    {\( y = (\e^{x} - 1)^{3}  \)（\(x \in \R\)）}
\end{problem}

\begin{answer}
D
\end{answer}

\begin{solution}
由 \(y=\ln(\sqrt[3]{x}+1)\) 得 \(\e^y=\sqrt[3]{x}+1\)，所以
\(x=(\e^y-1)^3\). 交换 \(x\)，\(y\) 后，反函数为
\(y=(\e^x-1)^3\). 原函数的值域为 \(\R\)，故反函数的定义域为
\(x\in\R\)，故选 D.
\end{solution}

\begin{problem}
已知 \(\bs{a}\)，\(\bs{b}\) 为单位向量，其夹角为 \(60^{\circ}\)，则 \((2\bs{a} - \bs{b}) \cdot \bs{b} =\)（\quad）

\choices
    {\(-1\)}
    {\(0\)}
    {\(1\)}
    {\(2\)}
\end{problem}

\begin{answer}
B
\end{answer}

\begin{solution}
由单位向量的夹角为 \(60^\circ\)，得
\(\bs{a}\cdot\bs{b}=|\bs{a}|\,|\bs{b}|\cos60^\circ=\frac{1}{2}\)，且
\(\bs{b}\cdot\bs{b}=1\). 因此
\((2\bs{a}-\bs{b})\cdot\bs{b}=2(\bs{a}\cdot\bs{b})-\bs{b}\cdot\bs{b}=1-1=0\)，故选 B.
\end{solution}

\begin{problem}
有\(6\text{ 名}\)男医生，\(5\text{ 名}\)女医生，从中选出\(2\text{ 名}\)男医生，\(1\text{ 名}\)女医生组成一个医疗小组，则不同的选法共有（\quad）

\choices
    {\(60\text{ 种}\)}
    {\(70\text{ 种}\)}
    {\(75\text{ 种}\)}
    {\(150\text{ 种}\)}
\end{problem}

\begin{answer}
C
\end{answer}

\begin{solution}
从 \(6\) 名男医生中选 \(2\) 名，从 \(5\) 名女医生中选 \(1\) 名，故不同选法共有
\(\mathrm C_6^2\mathrm C_5^1=15\times5=75\) 种，故选 C.
\end{solution}

\begin{problem}
设等比数列 \(\{a_{n}\}\) 的前 \(n\) 项和为 \(S_{n}\)，若 \(S_{2} = 3\)，\(S_{4} = 15\)，则 \(S_{6} =\)（\quad）

\choices
    {\(31\)}
    {\(32\)}
    {\(63\)}
    {\(64\)}
\end{problem}

\begin{answer}
C
\end{answer}

\begin{solution}
设等比数列的公比为 \(q\). 若 \(q=1\)，则 \(S_4=2S_2=6\)，与 \(S_4=15\) 矛盾，故 \(q\ne1\). 将前 \(4\) 项和分成前两项与后两项，得
\(S_4=S_2+q^2S_2=S_2(1+q^2)\). 因而
\(15=3(1+q^2)\)，即 \(q^2=4\). 再将前 \(6\) 项分成三个各含两项的部分，有
\[
S_6=S_2+q^2S_2+q^4S_2=3(1+4+16)=63
\]
故选 C.
\end{solution}

\begin{problem}
已知椭圆 \(C: \frac{x^2}{a^2} + \frac{y^2}{b^2} = 1 \)（\(a > b > 0\)）的左、右焦点分别为 \(F_1\)，\(F_2\)，离心率为 \(\frac{\sqrt{3}}{3}\)，过 \(F_2\) 的直线 \(l\) 交 \(C\) 于 \(A\)，\(B\) 两点，若 \(\triangle AF_{1}B\) 的周长为 \(4\sqrt{3}\)，则 \(C\) 的方程为（\quad）

\choices
    {\(\frac{x^2}{3} +\frac{y^2}{2} = 1\)}
    {\(\frac{x^2}{3} +y^2 = 1\)}
    {\(\frac{x^2}{12} +\frac{y^2}{8} = 1\)}
    {\(\frac{x^2}{12} +\frac{y^2}{4} = 1\)}
\end{problem}

\begin{answer}
A
\end{answer}

\begin{solution}
设椭圆的半焦距为 \(c\). 由离心率得
\(\frac{c}{a}=\frac{\sqrt{3}}{3}\)，所以 \(c=\frac{a}{\sqrt{3}}\)，从而
\(b^2=a^2-c^2=\frac{2a^2}{3}\). 因为 \(A\)，\(B\) 在椭圆上，且 \(F_2\) 在线段 \(AB\) 上，所以
\[
AF_1+AF_2=2a,\qquad BF_1+BF_2=2a,\qquad AF_2+BF_2=AB
\]
于是 \(\triangle AF_1B\) 的周长为
\(AF_1+BF_1+AB=4a\). 由题设 \(4a=4\sqrt{3}\)，得 \(a=\sqrt{3}\)，
\(c=1\)，\(b^2=2\)，故椭圆方程为
\(\frac{x^2}{3}+\frac{y^2}{2}=1\)，故选 A.
\end{solution}

\begin{problem}
正四棱锥的顶点都在同一球面上，若该棱锥的高为 \(4\)，底面边长为 \(2\)，则该球的表面积为（\quad）

\choices
    {\(\frac{81\pi}{4}\)}
    {\(16\pi\)}
    {\(9\pi\)}
    {\(\frac{27\pi}{4}\)}
\end{problem}

\begin{answer}
A
\end{answer}

\begin{solution}
设底面中心为 \(O\)，球心为 \(S\)，并设 \(OS=t\)，其中 \(S\) 在棱锥的高所在直线上. 正方形底面边长为 \(2\)，故底面中心到任一底面顶点 \(U\) 的距离为
\(\sqrt{2}\). 设棱锥顶点为 \(V\)，则 \(OV=4\)，由 \(SV=SU\) 得
\[
(4-t)^2=t^2+(\sqrt{2})^2=t^2+2
\]
解得 \(t=\frac{7}{4}\). 因而球半径
\(R=4-\frac{7}{4}=\frac{9}{4}\)，球的表面积为
\(4\pi R^2=4\pi\left(\frac{9}{4}\right)^2=\frac{81\pi}{4}\)，故选 A.
\end{solution}

\begin{problem}
双曲线 \(C: \frac{x^2}{a^2} - \frac{y^2}{b^2} = 1\)（\(a > 0\)，\(b > 0\)）的离心率为 \(2\)，焦点到渐近线的距离为 \(\sqrt{3}\)，则 \(C\) 的焦距等于（\quad）

\choices
    {\(2\)}
    {\(2\sqrt{2}\)}
    {\(4\)}
    {\(4 \sqrt{2}\)}
\end{problem}

\begin{answer}
C
\end{answer}

\begin{solution}
设双曲线的半焦距为 \(c\). 由离心率 \(e=\frac{c}{a}=2\)，得 \(c=2a\)，
再由 \(c^2=a^2+b^2\) 得 \(b^2=3a^2\)，即 \(b=\sqrt{3}a\).
双曲线的渐近线为 \(y=\pm\frac{b}{a}x\). 取右焦点 \(F=(c,0)\)，其到渐近线
\(bx-ay=0\) 的距离为
\[
\frac{bc}{\sqrt{a^2+b^2}}=\frac{\sqrt{3}a\cdot2a}{2a}=\sqrt{3}a
\]
由题设 \(\sqrt{3}a=\sqrt{3}\)，得 \(a=1\)，所以 \(c=2\)，焦距为 \(2c=4\)，故选 C.
\end{solution}

\begin{problem}
奇函数 \( f(x) \) 的定义域为 \( \R \)，若 \( f(x + 2) \) 为偶函数，\( f(1) = 1 \)，则 \( f(8) + f(9) = \)（\quad）

\choices
    {\(-2\)}
    {\(-1\)}
    {\(0\)}
    {\(1\)}
\end{problem}

\begin{answer}
D
\end{answer}

\begin{solution}
令 \(g(x)=f(x+2)\). 因为 \(g\) 为偶函数，所以
\(f(x+2)=f(-x+2)\). 令 \(t=x+2\)，得
\(f(t)=f(4-t)\). 由奇函数性质 \(f(0)=0\)，于是
\(f(4)=f(0)=0\). 又
\(f(5)=f(-1)=-f(1)=-1\)，因此
\(f(8)=f(-4)=-f(4)=0\)，
\(f(9)=f(-5)=-f(5)=1\). 所以 \(f(8)+f(9)=1\)，故选 D.
\end{solution}

\section{填空题}

\begin{problem}
\((x - 2)^{6}\) 的展开式中 \(x^{3}\) 的系数为\fillinblank{}.
\end{problem}

\begin{answer}
\(-160\)
\end{answer}

\begin{solution}
在二项式展开式中，含 \(x^3\) 的项取 \(x^3\) 与 \((-2)^3\)，故其系数为
\(\mathrm C_6^3(-2)^3=20\times(-8)=-160\).
\end{solution}

\begin{problem}
函数 \( y = \cos 2x + 2\sin x \) 的最大值为\fillinblank{}.
\end{problem}

\begin{answer}
\(\frac{3}{2}\)
\end{answer}

\begin{solution}
令 \(t=\sin x\)，则 \(t\in[-1,1]\)，且
\[
y=\cos2x+2\sin x=1-2\sin^2x+2\sin x
=\frac{3}{2}-2\left(t-\frac{1}{2}\right)^2\leq\frac{3}{2}
\]
当 \(t=\frac{1}{2}\) 时等号成立，所以最大值为 \(\frac{3}{2}\).
\end{solution}

\begin{problem}
设 \(x\)，\(y\) 满足约束条件 \(
\begin{cases}
x - y \geqslant 0, \\
x + 2y \leqslant 3, \\
x - 2y \leqslant 1
\end{cases}
\) 则 \(z = x + 4y\) 的最大值为\fillinblank{}.
\end{problem}

\begin{answer}
\(5\)
\end{answer}

\begin{solution}
由约束条件得可行域的三个顶点分别为
\[
(1,1),\qquad (-1,-1),\qquad \left(2,\frac{1}{2}\right)
\]
它们分别是直线 \(x-y=0\) 与 \(x+2y=3\)、\(x-2y=1\)，以及后两条直线的交点.
在三个顶点处，\(z=x+4y\) 的值分别为
\(5\)，\(-5\)，\(4\). 因此 \(z\) 的最大值为 \(5\).
\end{solution}

\begin{problem}
直线 \(l_{1}\) 和 \(l_{2}\) 是圆 \(x^{2} + y^{2} = 2\) 的两条切线，若 \(l_{1}\) 与 \(l_{2}\) 的交点为 \((1,3)\)，则 \(l_{1}\) 与 \(l_{2}\) 的夹角的正切值等于\fillinblank{}.
\end{problem}

\begin{answer}
\(\frac{4}{3}\)
\end{answer}

\begin{solution}
设 \(P=(1,3)\)，圆心为 \(O=(0,0)\)，切点为 \(T\). 由 \(PT\) 为切线，得
\(OT\perp PT\). 其中 \(OP=\sqrt{1^2+3^2}=\sqrt{10}\)，
\(OT=\sqrt{2}\)，所以
\(PT=\sqrt{OP^2-OT^2}=\sqrt{8}=2\sqrt{2}\).
设 \(OP\) 与其中一条切线的夹角为 \(\theta\)，则
\(\tan\theta=\frac{OT}{PT}=\frac{1}{2}\). 两条切线关于 \(OP\) 对称，故两直线的锐角为 \(2\theta\)，于是
\[
\tan(2\theta)=\frac{2\tan\theta}{1-\tan^2\theta}
=\frac{2\times\frac12}{1-\frac14}=\frac{4}{3}
\]
\end{solution}

\section{解答题}

\begin{problem}
数列 \( \{a_{n}\} \) 满足 \(a_{1}=1\)，\(a_{2}=2\)，\(a_{n+2}=2a_{n+1}-a_{n}+2.\)

\begin{enumerate}
\item 设 \( b_{n} = a_{n+1} - a_{n} \)，证明 \(\{b_{n}\}\) 是等差数列；
\item 求数列 \( \{a_{n}\} \) 的通项公式.
\end{enumerate}
\end{problem}

\begin{solution}
\begin{enumerate}
\item 由 \(b_n=a_{n+1}-a_n\)，得
\[
b_{n+1}=a_{n+2}-a_{n+1}=a_{n+1}-a_n+2=b_n+2
\]
又 \(b_1=a_2-a_1=1\)，所以 \(\{b_n\}\) 是首项为 \(1\)、公差为 \(2\) 的等差数列，从而
\(b_n=1+2(n-1)=2n-1\).
\item 由 \(a_n=a_1+\sum_{k=1}^{n-1}b_k\)，得
\[
a_n=1+\sum_{k=1}^{n-1}(2k-1)=1+(n-1)^2=n^2-2n+2
\]
\end{enumerate}
\end{solution}

\begin{problem}
\(\triangle ABC\) 的内角 \(A\)，\(B\)，\(C\) 的对边分别为 \(a\)，\(b\)，\(c\)，已知 \(3a\cos C = 2c\cos A\)，\(\tan A = \frac{1}{3}\)，求 \(B\).
\end{problem}

\begin{solution}
由正弦定理 \(\frac{a}{c}=\frac{\sin A}{\sin C}\)，将已知条件化为
\[
3\sin A\cos C=2\sin C\cos A
\]
因 \(\tan A=\frac13\)，且 \(A\in(0,\pi)\)，可知 \(A\) 为锐角，故 \(\cos A\ne0\). 若 \(\cos C=0\)，上式将推出 \(\cos A=0\)，矛盾，所以 \(\cos C\ne0\). 代入 \(\tan A=\frac13\)，得
\(\cos C=2\sin C\)，从而 \(\tan C=\frac12\)，且 \(C\) 也为锐角.
因此
\[
\tan(A+C)=\frac{\tan A+\tan C}{1-\tan A\tan C}
=\frac{\frac13+\frac12}{1-\frac16}=1
\]
又 \(A+C\in(0,\pi)\)，故 \(A+C=\frac{\pi}{4}\)，所以
\(B=\pi-(A+C)=\frac{3\pi}{4}\).
\end{solution}

\begin{problem}
如图，三棱柱 \(ABC - A_{1}B_{1}C_{1}\) 中，点 \(A_{1}\) 在平面 \(ABC\) 内的射影 \(D\) 在 \(AC\) 上，\(\angle ACB = 90^{\circ}\)，\(BC = 1\)，\(AC = CC_{1} = 2\).

\begin{enumerate}
\item 证明： \(AC_{1} \perp A_{1}B\)；
\item 设直线 \(AA_{1}\) 与平面 \(BCC_{1}B_{1}\) 的距离为 \( \sqrt{3} \)，求二面角 \( A_{1}-AB-C \) 的大小.
\end{enumerate}

\bitmapfigure[width=0.35\linewidth]{img/2014/syllabus_liberal/q19_fig1.png}
\end{problem}

\begin{solution}
\begin{enumerate}
\item 建立直角坐标系，令 \(C=(0,0,0)\)，\(A=(2,0,0)\)，\(B=(0,1,0)\)，并设
\(\overrightarrow{CC_1}=(u,0,h)\)，其中 \(h>0\). 这是因为 \(A_1\) 在底面 \(ABC\) 上的射影 \(D\) 位于 \(AC\) 上. 由棱柱的平移关系，
\[
A_1=(2+u,0,h),\qquad C_1=(u,0,h),\qquad u^2+h^2=CC_1^2=4
\]
于是
\[
\overrightarrow{AC_1}=(u-2,0,h),\qquad
\overrightarrow{A_1B}=(-2-u,1,-h)
\]
从而
\(\overrightarrow{AC_1}\cdot\overrightarrow{A_1B}=4-u^2-h^2=0\)，故 \(AC_1\perp A_1B\).
\item 平面 \(BCC_1B_1\) 的一个法向量为
\(\bs{n}=(h,0,-u)\)，而
\(\overrightarrow{AA_1}=(u,0,h)\)，所以
\(\bs{n}\cdot\overrightarrow{AA_1}=0\)，即直线 \(AA_1\) 平行于平面 \(BCC_1B_1\). 该平面过原点，方程为 \(hx-uz=0\)，故
\[
d(AA_1,\text{平面 }BCC_1B_1)
=\frac{|2h|}{\sqrt{h^2+u^2}}=h
\]
由题设 \(h=\sqrt3\)，结合 \(u^2+h^2=4\) 得 \(u^2=1\). 因 \(D=(2+u,0,0)\) 在线段 \(AC\) 上，故 \(u=-1\)，于是 \(A_1=(1,0,\sqrt3)\).
取
\(\overrightarrow{AB}=(-2,1,0)\)，\(\overrightarrow{AA_1}=(-1,0,\sqrt3)\)，则平面 \(A_1AB\) 的一个法向量为
\[
\bs{n}_1=\overrightarrow{AB}\times\overrightarrow{AA_1}=(\sqrt3,2\sqrt3,1)
\]
平面 \(ABC\) 的法向量可取 \(\bs{n}_0=(0,0,1)\). 设二面角 \(A_1-AB-C\) 为 \(\theta\)，则
\[
\cos\theta=\frac{|\bs{n}_0\cdot\bs{n}_1|}{|\bs{n}_0|\,|\bs{n}_1|}
=\frac{1}{\sqrt{3+12+1}}=\frac14
\]
由于 \(\theta\) 为锐角，\(\tan\theta=\sqrt{15}\)，所以 \(\theta=\arctan\sqrt{15}\).
\end{enumerate}
\end{solution}

\begin{problem}
设每个工作日甲，乙，丙，丁\(4\text{ 人}\)需使用某种设备的概率分别为 \(0.6\)，\(0.5\)，\(0.5\)，\(0.4\)，各人是否需使用设备相互独立.

\begin{enumerate}
\item 求同一工作日至少 \(3\) 人需使用设备的概率；
\item 实验室计划购买 \(k\) 台设备供甲，乙，丙，丁使用，若要求“同一工作日需使用设备的人数大于 \(k\)”的概率小于 \(0.1\)，求 \(k\) 的最小值.
\end{enumerate}
\end{problem}

\begin{solution}
\begin{enumerate}
\item 设甲、乙、丙、丁需使用设备分别为事件 \(A\)，\(B\)，\(C\)，\(D\). 由独立性，恰有 \(3\) 人需使用设备的概率为
\[
0.4\times0.5\times0.5\times0.4
+0.6\times0.5\times0.5\times0.4
+0.6\times0.5\times0.5\times0.4
+0.6\times0.5\times0.5\times0.6
=0.25
\]
四人都需使用设备的概率为
\(0.6\times0.5\times0.5\times0.4=0.06\). 因而同一工作日至少 \(3\) 人需使用设备的概率为 \(0.25+0.06=0.31\).
\item 设同一工作日需使用设备的人数为 \(X\). 当 \(k=2\) 时，
\(P(X>2)=P(X\geq3)=0.31>0.1\)，不满足要求；当 \(k=3\) 时，
\(P(X>3)=P(X=4)=0.06<0.1\). 由于 \(P(X>k)\) 随 \(k\) 增大而不增，故 \(k\) 的最小值为 \(3\).
\end{enumerate}
\end{solution}

\begin{problem}
函数 \( f(x) = ax^3 + 3x^2 + 3x \) （\( a \neq 0 \)）.

\begin{enumerate}
\item 讨论函数 \( f(x) \) 的单调性；
\item 若函数 \( f(x) \) 在区间\((1,2)\)是增函数，求 \( a \) 的取值范围.
\end{enumerate}
\end{problem}

\begin{solution}
\begin{enumerate}
\item \(f'(x)=3(ax^2+2x+1)\). 令 \(g(x)=ax^2+2x+1\)，其判别式为
\(\Delta=4(1-a)\).

当 \(a\geq1\) 时，\(g(x)\geq0\) 对一切 \(x\in\R\) 成立，所以 \(f(x)\) 在 \(\R\) 上是增函数.

当 \(0<a<1\) 时，令
\[
x_1=\frac{-1+\sqrt{1-a}}{a},\qquad x_2=\frac{-1-\sqrt{1-a}}{a}
\]
则 \(x_2<x_1\)，且 \(g(x)\) 在 \((-\infty,x_2)\) 和 \((x_1,+\infty)\) 上为正，在 \((x_2,x_1)\) 上为负. 因此 \(f(x)\) 在 \((-\infty,x_2)\) 和 \((x_1,+\infty)\) 上是增函数，在 \((x_2,x_1)\) 上是减函数.

当 \(a<0\) 时，仍令
\[
x_1=\frac{-1+\sqrt{1-a}}{a},\qquad x_2=\frac{-1-\sqrt{1-a}}{a}
\]
此时 \(x_1<x_2\)，且 \(g(x)\) 在 \((-\infty,x_1)\) 和 \((x_2,+\infty)\) 上为负，在 \((x_1,x_2)\) 上为正. 因此 \(f(x)\) 在 \((-\infty,x_1)\) 和 \((x_2,+\infty)\) 上是减函数，在 \((x_1,x_2)\) 上是增函数.
\item 要使 \(f(x)\) 在 \((1,2)\) 上为增函数，只需考察 \(g(x)\) 在该区间的符号. 当 \(a>0\) 时，\(g(x)>0\) 对 \(x\in(1,2)\) 恒成立. 当 \(a<0\) 时，\(g(x)\) 是开口向下的二次函数，在闭区间 \([1,2]\) 上的最小值取在端点，因此需有
\(g(1)=a+3\geq0\)，\(g(2)=4a+5\geq0\). 这两个条件等价于 \(a\geq-\frac54\)，与 \(a<0\) 合并得 \(a\in[-\frac54,0)\). 结合 \(a>0\) 的情形，得
\(a\in[-\frac54,0)\cup(0,+\infty)\).
\end{enumerate}
\end{solution}

\section{选考题：解答题}

\begin{problem}
已知抛物线 \(C: y^{2} = 2px \) （\(p > 0\)）的焦点为 \(F\)，直线 \(y = 4\) 与 \(y\) 轴的交点为 \(P\)，与 \(C\) 的交点为 \(Q\)，且 \(|QF| = \frac{5}{4} |PQ|\).

\begin{enumerate}
\item 求 \(C\) 的方程；
\item 过 \(F\) 的直线 \(l\) 与 \(C\) 相交于 \(A\)，\(B\) 两点，若 \(AB\) 的垂直平分线 \(l'\) 与 \(C\) 相交于 \(M\)，\(N\) 两点，且 \(A\)，\(M\)，\(B\)，\(N\) 四点在同一圆上，求 \(l\) 的方程.
\end{enumerate}
\end{problem}

\begin{solution}
\begin{enumerate}
\item 由 \(y=4\) 与抛物线 \(y^2=2px\) 联立，得
\(Q=(\frac{8}{p},4)\)，而 \(P=(0,4)\)，所以
\(|PQ|=\frac{8}{p}\). 抛物线的焦点为 \(F=(\frac p2,0)\)，于是
\[
|QF|^2=\left(\frac{8}{p}-\frac p2\right)^2+4^2
\]
由题设 \(|QF|=\frac54|PQ|=\frac{10}{p}\)，两边平方并乘以 \(p^2\)，得
\[
\left(8-\frac{p^2}{2}\right)^2+16p^2=100
\]
即 \(p^4+32p^2-144=0\). 因 \(p>0\)，解得 \(p^2=4\)，所以 \(p=2\). 因此
\(C\) 的方程为 \(y^2=4x\).
\item 此时 \(F=(1,0)\). 用抛物线的参数表示
\(T(r)=(r^2,2r)\). 若 \(l\) 为竖直线 \(x=1\)，则其弦 \(AB\) 的垂直平分线为 \(y=0\)，而 \(y=0\) 与抛物线只有一个交点，不符合 \(l'\) 与 \(C\) 相交于 \(M\)，\(N\) 两点的条件. 水平线 \(y=0\) 与抛物线也只有一个交点，不符合 \(l\) 与 \(C\) 相交于 \(A\)，\(B\) 两点的条件. 因此 \(l\) 可设为
\(l:y=t(x-1)\)，其中 \(t\ne0\). 交点 \(A=T(u)\)，\(B=T(v)\) 的参数满足
\[
2r=t(r^2-1)
\]
所以 \(u+v=\frac{2}{t}\)，\(uv=-1\). 记 \(s=u+v=\frac{2}{t}\). 则弦 \(AB\) 的中点为
\[
H=\left(\frac{u^2+v^2}{2},u+v\right)
=\left(\frac{s^2+2}{2},s\right)
\]
而弦 \(AB\) 的斜率为 \(t=\frac{2}{s}\). 因此其垂直平分线 \(l'\) 的方程为
\[
y-s=-\frac{s}{2}\left(x-\frac{s^2+2}{2}\right)
\]
若 \(M=T(m)\)，\(N=T(n)\)，将 \(T(r)=(r^2,2r)\) 代入 \(l'\)，得到
\[
2sr^2+8r-s(s^2+6)=0
\]
故 \(m+n=-\frac{4}{s}\).

再考虑圆与抛物线的交点参数. 任意圆可写为
\(x^2+y^2+\lambda x+\mu y+\nu=0\)，代入 \(T(r)\) 后得到
\[
r^4+(4+\lambda)r^2+2\mu r+\nu=0
\]
其中没有 \(r^3\) 项. 因而四点 \(A\)，\(M\)，\(B\)，\(N\) 共圆的必要条件为四个参数之和为 \(0\)，即
\[
u+v+m+n=s-\frac4s=0
\]
所以 \(s^2=4\)，从而 \(s=\pm2\)，\(t=\frac2s=\pm1\). 反过来，当 \(s=\pm2\) 时，四个参数的和为 \(0\)，其根式多项式没有三次项，因而可写成上述圆与抛物线联立所得的形式，故四点确实共圆. 因此
\[
l:y=x-1\quad\text{或}\quad y=-x+1
\]
\end{enumerate}
\end{solution}
